\documentclass[11pt, a4paper]{article}
\usepackage[utf8]{inputenc}
\usepackage{geometry}
\usepackage{amsmath}
\usepackage{graphicx}
\usepackage{hyperref}
\usepackage{natbib}
\usepackage{titlesec}
\usepackage{setspace}

% Geometry setup
\geometry{top=2.5cm, bottom=2.5cm, left=2.5cm, right=2.5cm}

% Formatting
\titleformat{\section}{\large\bfseries}{\thesection}{1em}{}
\onehalfspacing

% Colors for hyperlinks
\hypersetup{
    colorlinks=true,
    linkcolor=blue,
    filecolor=magenta,
    urlcolor=cyan,
    citecolor=blue,
}

\title{\textbf{Inverting the Lithospheric Stress Field of Central Iran: \\ A Physics-Informed Deep Learning Approach using Sparse Kinematic Data}}
\author{\textbf{Research Proposal}}
\date{\today}

\begin{document}

\maketitle

\begin{abstract}
Accurate estimation of the lithospheric stress field is critical for earthquake forecasting but remains challenging due to sparse \textit{in-situ} measurements. Traditional Finite Element Models (FEMs) require complex meshing and explicit boundary conditions that are often unknown. We propose a novel \textbf{Physics-Informed Neural Network (PINN)} framework to invert for the 3D continuous stress tensor of Central Iran. Unlike previous studies that assume purely elastic rheology, our model incorporates \textbf{Maxwell Viscoelasticity} to capture time-dependent stress relaxation. Furthermore, we introduce a method to constrain the network using only \textbf{sparse GPS azimuths} (orientations), solving a highly non-linear inverse problem. This approach promises a mesh-free, physics-consistent stress map that outperforms statistical baselines (XGBoost) and standard Coulomb Stress Transfer (CST) in forecasting seismicity.
\end{abstract}

\section{Scientific Rationale \& Gap Analysis}
Current state-of-the-art methods generally fall into two categories:
\begin{enumerate}
    \item \textbf{Data-Driven (e.g., XGBoost)}: Excellent at finding correlations but lacks physical constraints, leading to poor generalization in data-sparse regions.
    \item \textbf{Classical Physics (e.g., FEM/CST)}: Physically rigorous but limited by mesh quality and the need for unknown boundary conditions.
\end{enumerate}

\subsection*{The Gap}
Recent works (e.g., \textit{Zhang et al., Sci. Rep. 2022}) have demonstrated PINNs for stress inversion but rely on simplified \textbf{Linear Elastic} assumptions and dense synthetic data. No study has yet applied \textbf{Viscoelastic PINNs} to invert for the stress field of the complex \textbf{Central Iranian collision zone} using real, sparse \textbf{GPS strain azimuths}.

\section{Methodology}

\subsection{Physics-Informed Architecture}
We define a deep neural network $\mathcal{N}_\theta(x,y,z,t)$ that approximates the stress tensor $\sigma_{ij}$ and velocity field $v_i$. The Loss Function $\mathcal{L}$ minimizes the residuals of the momentum balance and constitutive equations:
\begin{equation}
    \mathcal{L}_{total} = \mathcal{L}_{metric} + \lambda_1 \mathcal{L}_{PDE} + \lambda_2 \mathcal{L}_{BC}
\end{equation}
where:
\begin{itemize}
    \item $\mathcal{L}_{PDE}$ (Momentum): $\nabla \cdot \sigma + \rho g = 0$
    \item $\mathcal{L}_{PDE}$ (Constitutive): $\dot{\epsilon}_{ij} = \frac{1}{2\mu}\dot{\sigma}_{ij} + \frac{1}{2\eta}\sigma_{ij}$ (Maxwell Viscoelasticity)
\end{itemize}

\subsection{Data Integration (The "Inverse" Problem)}
We utilize the provided \texttt{kinematic\_data} (GPS strain rates). Crucially, these data provide \textbf{orientations} (azimuths) rather than absolute stress magnitudes. We enforce this constraint by penalizing the misalignment between the predicted principal stress direction and the observed GPS strain azimuth:
\begin{equation}
    \mathcal{L}_{data} = \sum_{k} (1 - \cos(\theta_{pred}^{(k)} - \theta_{obs}^{(k)}))^2
\end{equation}

\subsection{Execution Roadmap}
\begin{enumerate}
    \item \textbf{Phase I (Validation)}: Implement a 2D Plane Stress model to validate the formulation against a synthetic "toy" fault interaction.
    \item \textbf{Phase II (Main)}: Scale to full 3D Lithospheric model using the Central Iran velocity structure (\texttt{Pwave.3D.txt}) as a proxy for heterogeneous material properties ($\mu(x)$).
\end{enumerate}

\section{Validation \& Benchmarking}
To satisfy the rigor required by high-impact journals, we will benchmark against:
\begin{enumerate}
    \item \textbf{Coulomb Stress Transfer (CST)}: The standard physics-based predictor.
    \item \textbf{XGBoost (Spatial Susceptibility)}: The standard ML baseline.
\end{enumerate}
We will demonstrate that the PINN model achieves superior log-likelihood in \textbf{hindcasting} historical seismicity (2006-2023) compared to these baselines.

\section{Expected Impact}
This study will provide the first-ever continuous, physics-consistent 3D stress map of Central Iran. By demonstrating the efficacy of Viscoelastic PINNs in data-sparse regions, we establish a new framework for earthquake potential imaging applicable to collision zones worldwide.

\end{document}
